\chapter*{译者说明}
\addcontentsline{toc}{chapter}{译者说明}

本书是 Johannes Anschütz《Lecture notes for an introductory course on the Langlands program》（2020 年 7 月 22 日修订稿，103 页）的中文译注本，其中第 7、8、10 讲由 Ben Heuer 撰写，标题中已分别标明。

\section*{底本与译名}

底本为作者在波恩大学开设的朗兰兹纲领入门课程的讲义，共 12 讲：前六讲从模形式出发，经 adèle 化、Hecke 代数与强重数一，把新形式组织成自守表示，并陈述 $\mathrm{GL}_2$ 上的 Langlands 互反；第 7、8 讲（Heuer）讲模曲线作为椭圆曲线模空间、Eichler--Shimura 关系；第 9 讲用上同调实现附着于新形式的 Galois 表示；第 10 讲（Heuer）用模 $p$ 同余与 Deligne--Serre 提升处理权 $1$ 形式；第 11、12 讲把整套陈述推广到一般约化群（$L$-群、Satake 参数、$L$-参数、Buzzard--Gee 猜想），并以 Shimura 簇收尾。

译文尽量保留讲义的条目层级与标签（「上次」「要点」「事实」「注」「略述」），文献引用（如 \texttt{[GH19, Lemma 3.7.1]}）与定义、定理、公式的编号一律原样保留，不重编。

\section*{体例}

\begin{itemize}
\item 页边不显示原文页码；每页正文前在源码中以注释形式保留页标记，供校对核对，不进入成品。
\item 译者添加的解释以脚注形式出现，标有「【译注】」者均有可追溯来源（见脚注内的\emph{来源}），标有「【译者按】」者为本译本的补充说明、无外部依据。
\item 原文中的笔误（例如 \S 5 中 $\Gamma_0(m)$ 的 $\Gamma_0$ 条件、\S 4 中 $Y$ 的矩阵分解、\S 11 中 Casimir 元的一段中间计算等）均在译注中说明原文写法与校正依据，正文按校正后的含义翻译。
\item 术语与同主题前作《Langlands 纲领导论》（SY02-09）对齐；少数与《朗兰兹纲领导论（全书中文译注）》（SY02-24）不同的译法，在首次出现处以括号附英文原词。
\end{itemize}

\section*{致谢与声明}

感谢原作者 Johannes Anschütz 与 Ben Heuer 写出并公开这份讲义。本译注本为非商业学术译本，底本版权归原作者所有；若原作者不愿本译本流传，请告知，本库将立即撤下。译文中一切错误由本译本负责。
